\chapter{\textcolor{Salmon}{Dashboards als potenzielle Lösung von Effizienz- und Übersichtlichkeitsproblemen in Legacy-Test-Systemen}}
\label{ch:Problemstellung}
\textcolor{blue}{\textit{einleitende Worte}}\\
ordnet ein: gibt einen Überblick von den Grundlagen zur Problemstellung\\
-> bis zur Vorstellung der Lösung (Methodik)

% Struktur und Bedeutung von Kapitel 3:     %// TODO Prof
%Kapitel 3 scheint aktuell hauptsächlich aus eigenen Hinweisen für dich selbst zu bestehen. Bitte beachte: Dieses Kapitel ist inhaltlich sehr wichtig, da es den Rahmen für die folgenden Kapitel bildet. Eine klare, strukturierte Darstellung in Kapitel 3 hilft der Leserschaft, den weiteren Verlauf besser einzuordnen.

\section{\textcolor{Salmon}{Herausforderungen in der Erstellung von Dashboards und der Arbeit mit Legacy- und Code-Testing-Systemen}}
\label{ch:Problem:Gesamt}

Veraltete Systeme stellen Unternehmen zunehmend vor Herausforderungen. Die Systeme sind groß, umfrangreich, gewachsen und wenig getestet. Sie bestehen aus vielen verschiedenen Technologien und sind schwierig zu warten. Die Fachkräfte, welche die Systeme aufgebaut und Jahrelang betreut habe stehen kurz vor der Rente und der Erhalt der Altsysteme ist langfristig nicht der beste Weg.

Zuvor wurde in \autoref{ch:SdF:Legacy:sec:Methoden} primär darauf eingegangen, wie Legacy-Code zu modernisieren ist. Dabei resultiert der scheinabr beste Ansatz, die Systeme möglichst frühzeitig, Stückweise mit dem Technikastand mitzuführen bzw. die Systeme sukzessiv nachzurüsten. Vor allem umfangreiche und gepflegte Testumgebungen vereinfachen die Modernisierung, indem das gewollte Systemverhalten dauerhaft nachgeprüft werden kann.

Dabei ist dieser Testansatz nicht neu. Unternehmen, die sich über die jahre Testsysteme und Strukturen aufgebaut haben, können auch hier vor dem Problem stehen, dass diese Systeme sehr umfangreich gewachsen sind (s. \autoref{ch:SdF:Legacy:sec:Legacy_ungetestet}). Mit zunehmenden Maschinen, Versionen, Komponenten und Testfällen, vermehren sich die zu überwachenden Ergebnisse rasant.

Trotz zusammenfassender Ergebnisübersichten ist es schwierig, den Überblick zu behalten und garantieren zu können, alle Werte überprüft zu haben. Zudem wird zunehmend mehr Zeit zu diesem manuellen Monitoring benötigt. Durch fehlende Motivation, abnehmende Aufmerksamkeitsspanne sowie Routine, sind Fehler im Prozess, wie übersehende Testergebnisse oder eine fehlerhafte Interpretation dieser möglich. Diese Arbeit ist ohne Automation somit höchst ineffizient.

Werden neue Systeme zur Ergebnisanzeige, vor allem Dashboards, umgesetzt, muss das System zuerst analysiert werden, um den gesamten Umfang zu erfassen. Damit ein solches neues System auch genutzt wird, muss eine Vollständigkeit garantiert werden können, um nicht zwischen den Systemen im Prozess wechseln zu müssen.

Eine weitere Schwierigkeit besteht darin, dass gewachsene und selbstgebaute Systeme häufig nicht kompatibel mit modernen Standards sind. Hierfür muss das System entweder von Grund auf neu geschrieben werden oder entsprechende Adapter implementiert werden.

\section{\textcolor{Salmon}{Fokussierung auf die Problemstellung der Effizienz und Übersichtlichkeit verteilter Testsysteme}}
1-2 Probleme fokussieren davon (von a)\\
-> unübersichtliche Umgebung - Fokussierung auf die Problemstellung der Übersichtlichkeit verteilter Testsysteme\\
-> auf eigenes Problem anpassen\\

%von a) ableiten, was mein Problem darin ist (halbe Seite)

\underline{Usability:}
\begin{itemize}
    \item In Form von Effektivität, Effizienz und Zufriedenheit
    \item Einfluss von Übersichtlichkeit -> Verteilung der Daten, des Systems
\end{itemize}

aus den Probelmen eine Sichtweise rausstellen
--> daraus stellt sich die Fragestellung xy

Hauptfrage - die aufsplitten in Teilfragem
--> Diskussion: Einzelfragestellung
--> Ergebnisse in ein Ergebnis zusammenfassen
\section{\textcolor{Salmon}{Beitrag moderner Dashboards zur Effizienzsteigerung und verbesserten Übersichtlichkeit von Legacy-Testsystemen}}\\

%Fragestellung von eigenem Problem ableiten, ohne eine Frage zu formulieren\\
%-> Subfragen bezüglich Ergebnissen, die eine Aussage zur Usability treffen\\

%kurz, nur Auflistung der Fragen\\
\noindent
\textbf{Titel}: Analyse der Steigerung der Usability durch Einführung eines Dashboards in den Legacy-Testpoolmonitoring-Prozess\\

\noindent
- aufschreiben, wie die Antworten die Thesen supporten oder nicht\\
-> die vorhergestellten Fragen\\
-> „was ist das erwartete Ergebnis“ (Regelwerk – es wird Ergebnis „turning blue“ erwartet wenn wir es richtig machen)\\
-> wie welche Antworten meine These unterstreichen\\

\noindent
- Wie passen die Hypothesen zu den Fragestellungen??\\

\noindent
- Wie ist der Stand?\\
- Worum geht es?\\
- Was ist das Ziel?\\
-> Forschungsfragen formulieren\\

\noindent
in Problemstellung:\\
- Übergang von SdF auf: was hat das mit meinem Thema zu tun\\
- wie lassen sich die Ansätze neu interpretieren auf das gegebene Problem

\noindent
Dashboard als nicht tiefe Lösung - ohne Eingriff in  das bestehende System\\
-> als erster Schritt zur Verbesserung der Usability\\
(Fazit: die die sich nicht mehr umgewöhnen wollen / können, müssen es mit diesem Ansatz nicht)

\noindent
\underline{aus Exposé}\\
Die Hauptfragestellung lautet: Welche Auswirkungen hat die Einführung eines Dashboards auf die Effizienz eines Legacy-Prozesses, untersucht am Beispiel des Testpoolmonitorings? \\
Dabei sollen folgende Teilfragen beantwortet werden:
\begin{itemize}
    \item Welche Herausforderungen bestehen, bei der Modernisierung von Legacy-Systemen?
    \item Welche Vorteile bringt ein Dashboard einem Legacy-System?
    \item Welche Rolle spielt die Akzeptanz im Nutzerumfeld, damit Modernisierungen sinnvoll sind?
\end{itemize}
\section{\textcolor{Plum}{Methodisches Vorgehen zur Entwicklung eines Dashboards als
moderne Anzeige der Testpoolergebnisse}}
...

\subsection{\textcolor{Plum}{Analyse des Ist-Zustandes der Testpool-Ergebnisanzeige zur Identifikation bestehender Anforderungen}}
...

\subsection{\textcolor{Plum}{Entwicklung eines Dashboards als moderne Ergebnisübersicht}}
...

\subsection{\textcolor{Plum}{Usability-Test zur Validierung der gestellten Anforderungen an das Dashboard}}
...

\subsection{\textcolor{Plum}{Nutzerbefragung zur Validierung der Benutzererfahrung im
Rahmen des Usability-Tests}}
...
